现在做数据的人，手上的特征数和样本量常常是同一个量级。基因表达是几千个基因对几百个样本，文本是几万维词表对几万篇文档，推荐系统的特征平台随手就能拼出几百上千维，而可用的标注样本未必更多。经典数理统计默认的场景是维数固定、样本量趋于无穷；这个单元处理的是另一侧：维数与样本量之比 \(c=p/n\) 不可忽略时会发生什么。

先看一个能自己跑一遍的现象。生成 \(1000\times1000\) 的标准正态数据矩阵，真协方差就是单位阵，一千个特征值全都等于 1。算样本协方差的特征值，得到的不是一千个贴着 1 的数，而是一片从 0 铺到 4 的连续分布。把 \(p\) 和 \(n\) 同时放大十倍、保持 \(c=1\)，直方图的形状纹丝不动。样本量增加解决不了它——这不是有限样本的误差，是高维极限下的确定形状。

这件事落到日常工作里就是：碎石图的肘部可能是噪声画出来的，PCA 报出的解释方差比系统性偏高，主成分方向在重采样时来回摆，而任何需要求协方差逆的方法（马氏距离、判别分析、组合优化）都会被压到零附近的最小特征值放大到失控。随机矩阵理论的用处，是把这些症状换成可计算的偏差。

\subsection{三节各自回答什么}

第一节确认现象并把它和随机波动区分开：谱的外扩是由 \(c\) 决定的系统性偏差，逐元素无偏挡不住它，只有 \(c\to0\) 才会消失。第二节给出精确形状——Marchenko--Pastur 律把``纯噪声该长什么样''写成一条闭式密度曲线，支撑区间 \([(1-\sqrt c)^2,(1+\sqrt c)^2]\) 的上界就成了一条可以直接用的信号阈值线，同时也要看清这条线在什么条件下不成立。第三节把偏差传到下游：BBP 相变说明信号强度低于 \(\sqrt c\) 时，估计出的主方向与真方向渐近正交；而 Johnson--Lindenstrauss 引理说明随机低维投影反倒能保住距离，所需维数只有 \(O(\log n/\varepsilon^2)\)，与原始维数无关。

第三节这个对照是本单元的落点。同一套随机矩阵的谱工具，一边警告高维下估计出来的方向不可信，一边说明高维下随手抽的方向很好用。差别不在维数本身，在于你是否需要用有限样本去估计一个高维方向。

\subsection{本章篇目}

以下篇目按概念依赖排列；若从具体问题进入，也可以先打开最接近当前困惑的一篇，再沿篇内链接回补。

\begin{itemize}
\tightlist
\item \hyperref[sec:11-Mathematical-Statistics/09-Random-Matrix-Theory/01]{``样本协方差在高维为何会失真''}；
\item \hyperref[sec:11-Mathematical-Statistics/09-Random-Matrix-Theory/02]{``Marchenko--Pastur 律给出谱的精确形状''}；
\item \hyperref[sec:11-Mathematical-Statistics/09-Random-Matrix-Theory/03]{``谱偏差怎样影响 PCA 与随机投影''}。
\end{itemize}

读完三节之后，再看 \hyperref[sec:03-Linear-Algebra/07-SVD-and-Data-Applications/04]{``从数据矩阵到 PCA''} 与 \hyperref[sec:13-Machine-Learning/06-Probabilistic-Models/02]{``LDA 与收缩：判别方向的几何与高维失效''}，那里的几何论证仍然成立，只是每个读数都得先按 \(c\) 折一次算。
